\chapter{Glossario}
\newpage
\begin{table}[H]
    \small
    \centering
    \caption{Glossario dei termini principali del sistema Wiki Collaborativo}
    \label{tab:glossario}
    \begin{tabularx}{\textwidth}{@{}l X@{}}
        \toprule
        \textbf{Termine} & \textbf{Definizione} \\
        \midrule
        \textbf{Amministratore} & Utente con privilegi elevati incaricato della gestione strutturale di corsi e argomenti. \\
        & \textit{Sinonimi: Admin.} \\
        \midrule
        \textbf{Appunto} & L'unità informativa principale del sistema, contenente testo in formato Markdown o plain text renderizzato in HTML. \\
        \midrule
        \textbf{Argomento} & Sotto-categoria tematica di un corso a cui vengono associati uno o più appunti. \\
        \midrule
        \textbf{Corso} & Categoria tematica a cui vengono associati uno o più argomenti. \\
        \midrule
        \textbf{CRUD} & Acronimo di Create, Read, Update, Delete; rappresenta le quattro operazioni di base gestite dalle API REST del sistema. \\
        \midrule
        \textbf{Markdown} & Linguaggio di markup leggero utilizzato dagli studenti per la formattazione del testo degli appunti. \\
        \midrule
        \textbf{MVP} & Pattern architetturale (Model-View-Presenter) che separa la logica di business dalla visualizzazione, come richiesto dai requisiti di progetto. \\
        \midrule
        \textbf{REST API} & Interfaccia programmabile stateless che permette la comunicazione tramite protocollo HTTP tra il frontend (JS) e il backend (PHP). \\
        \midrule
        \textbf{Studente} & Utente registrato che ha effettuato l'accesso e possiede i permessi per creare o modificare appunti e visualizzare la cronologia. \\
        & \textit{Sinonimi: Autore, Stud., StudAuth.} \\
        \midrule
        \textbf{Versione} & Stato storico di un appunto salvato in seguito a una modifica o a un ripristino. Consente di navigare nella cronologia dei contenuti e di recuperare checkpoint precedenti. \\
        & \textit{Sinonimi: Storico, Cronologia, Checkpoint.} \\
        \midrule
        \textbf{Visitatore} & Utente non registrato che non ha effettuato l'accesso. Può solo visualizzare e cercare gli appunti senza poterli creare o modificare. \\
        & \textit{Sinonimi: Visit.} \\
        \bottomrule
    \end{tabularx}
\end{table}